\section{导读：AI 下半场从模型能力走向系统能力}

本节先建立整期的坐标。姚顺雨在 2025 年 4 月发布《The Second Half》，认为 AI 主线程已经进入下半场。上半场的主线更像是模型能力：预训练、规模化、对话和通用能力；下半场则更像系统能力：Agent、任务环境、奖励机制、工具调用、交互界面和组织结构。访谈从个人经历讲起，但真正的核心是一个问题：当模型足够强以后，智能怎样进入世界、形成行动、组织和新的边界？

\begin{figure}[H]
\centering
\includegraphics[width=0.92\textwidth]{second-half-map.png}
\caption{AI 下半场地图：从模型能力走向 Agent、系统、交互和人类全局。自制概念图，依据 00:01:45--02:31:20 对谈内容整理。}
\end{figure}

\begin{knowledgebox}{读图：下半场不是“模型不重要”}
模型仍然是底座，但讨论重心从单个模型分数转向系统如何行动：任务是否真实、环境是否足够大、reward 是否可定义、interface 是否打开新机会、最终世界是单极还是多元。
\end{knowledgebox}

\begin{importantbox}{本期核心命题}
Agent 不是聊天模型的一个功能按钮，而是一种系统范式：它要在环境里观察、推理、行动、获得反馈，并通过工具、奖励和多智能体结构不断扩大可完成任务的边界。
\end{importantbox}

\begin{knowledgebox}{核心概念总表}
\begin{tabularx}{\textwidth}{p{0.20\textwidth}p{0.34\textwidth}X}
\toprule
概念 & 本期含义 & 学习时要避免的误解 \\
\midrule
Agent & 与环境交互并试图优化目标的系统 & 不是“聊天框加几个工具”这么简单。 \\
Environment & Agent 行动和获得反馈的世界 & 环境太小，能力就无法外推。 \\
Reward & 评价动作后果的信号 & reward 不等于人类真正想要的全部价值。 \\
Affordance & 环境给行动者提供的行动可能性 & code/API 是 AI 的手，不只是输出格式。 \\
Interface & 用户和智能系统交互的方式 & UI 会改变任务、数据、商业和研究方向。 \\
Different Bet & 与霸主不同的下注路线 & 不是换皮复制，而是改变边界条件。 \\
\bottomrule
\end{tabularx}
\end{knowledgebox}

\begin{knowledgebox}{本期阅读路线}
\begin{tabularx}{\textwidth}{p{0.22\textwidth}p{0.32\textwidth}X}
\toprule
路线 & 核心问题 & 对应章节 \\
\midrule
人 & 为什么语言是泛化工具，为什么姚顺雨押注 Agent？ & 语言与非共识。 \\
系统 & Agent 的任务、环境、方法线如何共同演化？ & 两个核心问题与三波兴衰。 \\
奖励 & Agent 如何拥有 reward、自我探索和 multi-agent 结构？ & reward、code affordance、泛化。 \\
边界 & Chatbot、interface、Super App 怎样决定新机会？ & 吞噬边界与交互方式。 \\
全局 & 单极与多元、OpenAI 与创业机会如何共存？ & 人类全局与 different bet。 \\
\bottomrule
\end{tabularx}
\end{knowledgebox}

\subsection{本章小结}

EP115 是 Agent 理论访谈，不只是 OpenAI 研究员访谈。它把语言、人、任务、环境、奖励、交互和平台竞争放在同一张图里，适合作为理解 Agent 下半场的核心笔记。

